\documentclass[12pt,a4paper]{report}
\usepackage[]{amsmath,amsthm,amssymb,amscd}
\usepackage[latin1]{inputenc}
%\usepackage{fullpage}

\topmargin 0pt
\advance \topmargin by -\headheight
\advance \topmargin by -\headsep
     
\textheight 246mm
     
\oddsidemargin 0pt
\evensidemargin \oddsidemargin
\marginparwidth 0.5in
     
\textwidth 159mm


%               Theorem Environments
\theoremstyle{definition}
\newtheorem{ex}{}
\newcommand{\pd}[1]{\frac{\partial}{\partial #1}} %\pd{x}
%               Subquestions numbered with letters
\renewcommand{\theenumi}{\textbf{\alph{enumi}}}

%               Maths definitions

\newcommand{\NN}{\ensuremath{\mathbb N}}
\newcommand{\ZZ}{\ensuremath{\mathbb Z}}
\newcommand{\CC}{\ensuremath{\mathbb C}}
\newcommand{\RR}{\ensuremath{\mathbb R}}
\newcommand{\TT}{\ensuremath{\mathbb T}}
\newcommand{\g}{\ensuremath{\frak{g}}}
\newcommand{\h}{\ensuremath{\frak{h}}}
\newcommand{\ft}{\ensuremath{\frak{t}}}
\newcommand{\cB}{\mathcal{B}}
\newcommand{\cN}{\mathcal{N}}
\newcommand{\cL}{\mathcal{L}}
\newcommand{\cD}{\mathcal{D}}
%\newcommand{\pd}[1]{\frac{\partial}{\partial #1}} %\pd{x}
%\newcommand{\pd}[1]{\frac{\partial}{\partial #1}} %\pd{x}

\begin{document}
\noindent
{\sffamily\large Marco Zambon 
\hfill   a discutir:   Miercoles, 16/05/2012\\Geometr\'ia III, curso 2011-12}
 

\noindent
\hrulefill{}\\
\begin{center}\bfseries\large\textsf{Hoja 6    \vspace{-2mm}}
\end{center}

\noindent
%\hrulefill{}\\
% \begin{ex} 
%Sea $F \colon M\to N$ una submersi\'on sobreyectiva entre variedades.
%Demuestra para todo $k\ge 0$ que la aplicaci\'on lineal (definida en clase)
%$$\Omega^k(N) \to  \Omega^k(M) , \alpha \mapsto F^* \alpha$$
%es inyectiva. 
% 
%\noindent{\textbf{Recuerda:}   una aplicaci\'on diferenciable $F$ es una submersi\'on si en todo $p\in M$, la derivada $F_*(p)$ es sobreyectiva.}
%\end{ex}

 
 \begin{ex} Sea $M=\RR^2$. 
 \begin{enumerate}
\item Existe $f\in C^{\infty}(M)$ tal que $df=x dy$?
\item Existe $f\in C^{\infty}(M)$ tal que $df=x dy+y dx$?
\end{enumerate}
\end{ex}
 
\begin{ex} \label{int}
Sea $M$ una variedad, $\alpha\in \Omega^1(M)$ una 1-forma, y 
$c \colon [a,b]\to M$ una curva diferenciable. Definimos  \emph{la integral de $\alpha$ a lo largo de $c$}:
$$\int_c\alpha:=\int_a^b \alpha_{c(t)}(\dot{c}(t)) dt.$$
Demuestra que, si $f$ es una funci\'on diferenciable en $M$, entonces
$$\int_{c} df=f(c(b))-f(c(a)).$$ 
\end{ex}



\begin{ex} 
Sea $(M,g)$ una variedad Riemanniana y $f \in C^{\infty}(M)$. El \emph{gradiente} de $f$ es el campo vectorial $grad(f)$ en $M$ definido por la propriedad:
$$g(grad(f),X)=df(X) \text{ para todo campo vectorial $X$ en $M$}.$$
Demuestra que:
\begin{enumerate}
\item Para cada $p\in M$ tal que $df_p \neq 0$   la funci\'on
$$\{Y\in T_pM : |Y|=1\}\to \RR, Y \mapsto df_p(Y)$$
tiene su valor m\'aximo en $\frac{grad(f)_p}{|{grad(f)_p}|}$. Aqu\'i utilizamos la notaci\'on $|Y|:=\sqrt{g_p(Y,Y)}$.
\item Sea $c \colon [a,b] \to M$ una curva integral de $grad(f)$. Entonces $$f(c(b))\ge f(c(a)).$$
\end{enumerate}


\noindent\textbf{Consejo para a):} utiliza la desigualdad de Cauchy-Schwarz.\\
\noindent\textbf{Consejo para b):} utiliza el ejercicio \ref{int}.
\end{ex}

\newpage
\begin{ex} 
Sea $$O(3)=\{A\in Mat(3,\RR):A^TA=Id\}.$$ Cada $A\in O(3)$ induce un difeomorfismo $\Phi_A \colon \RR^3 \to \RR^3, w \to A\cdot w$.
Adem\'as, 
sea $S^2:=\{(x,y,z)\in \RR^3:x^2+y^2+z^2=1\}$ la esfera. Demuestra:

\begin{enumerate}
\item Para cada $A\in O(3)$, $\Phi_A$ es una isometr\'ia de $\RR^3$ (con su metrica Riemanniana \'estandard $g$).
\item Para cada $A\in O(3)$, la aplicaci\'on $\Phi_A$ satisface $\Phi_A(S^2)=S^2$. Denotamos $\phi_A \colon S^2 \to S^2$ la aplicaci\'on inducida por restricci\'on.
\item\label{c} Para cada $A\in O(3)$, la aplicaci\'on $\phi_A$ es una isometr\'a de $S^2$ (con la metrica Riemannniana $\bar{g}$ inducida por la metrica est\'andard en $\RR^3$).
\item Deduce de \ref{c}. que, dados puntos $p,q\in S^2$, existe una isometr\'ia $\phi$ de $S^2$ con $\phi(p)=q$.
\item Deduce de \ref{c}. que el plano proyectivo $P_2(\RR)$ tiene una metrica Riemanniana tal que la proyecci\'on $\pi \colon S^2 \to P_2(\RR)$ es una isometr\'ia local.
\end{enumerate}

\noindent\textbf{Recuerda (para a):} Por las clases de calculo, como $\Phi_A$ es una aplicaci\'on lineal, tenemos $D_p \Phi_A=A$ en cada $p\in \RR^3$. 

\noindent\textbf{Recuerda (para d):} Las matrices de $O(3)$ son exactamente las matrices cuyas columnas forman una base ortonormal de $\RR^3$.

\noindent\textbf{Observaci\'on (para e):} $\pi$ es una    isometr\'ia local si  cada $p\in S^2$ tiene un entorno $U$ tal que $\pi|_U \colon U \to \pi(U)$ es una isometr\'ia.
\end{ex}

\begin{ex}
Considera el paraboloide $$P=\{(x_1,x_2,x_3)\in \RR^3:x_1^2+x_2^2=x_3\},$$ con la metrica Riemanniana inducida por $\RR^3$. Describe  la geod\'esica $\gamma$ que satisface $$\gamma(0)=(0,0,0)\;\;\;\;\;\;\;{\gamma}^{\prime}(0)=\pd{x_1}|_{(0,0,0)}.$$

\noindent\textbf{Consejo:} Se puede decir bastante sobre $\gamma$ sin hacer ninguna cuenta, utilizando el hecho que $\Phi \colon \RR^3 \to \RR^3, (x_1,x_2,x_3) \mapsto (x_1,-x_2,x_3)$ induce por restricci\'on  una isometr\'ia de $P$.
\end{ex} 
 

%\begin{ex}
%Sea $(M,g)$ una variedad Riemanniana compacta. Demuestra que existen $p_0,q_0\in M$ tal que para cualquier otro par de puntos $p,q\in M$, tenemos $$d(p_0,q_0)\ge d(p,q).$$ Demuestra que, si quitamos  la suposici\'on que 
%$M$ sea compacta, la afirmaci\'on es falsa.

%\noindent\textbf{Consejo:} Demuestra que la aplicaci\'on $M \times M \to \RR, (p,q)\to d(p,q)$ es continua, es decir, preimagenes de abiertos de $\RR$ son abiertos en $M\times M$. 
%\end{ex}    
 
  \end{document}
 %%%%%%% 
 
 

 \begin{ex}
\begin{enumerate}
\item 
\end{enumerate}
\end{ex}
 
 
 
 